% Generated draft: requires mathematical and source review.
\documentclass[11pt]{article}
\usepackage[margin=25mm]{geometry}
\usepackage{fontspec}
\setmainfont{Latin Modern Roman}
\usepackage{amsmath,amssymb,mathtools,unicode-math}
\setmathfont{Latin Modern Math}
\usepackage{enumitem,needspace,xurl,hyperref}
\setlength{\parindent}{0pt}
\setlength{\parskip}{5pt}
\setlength{\emergencystretch}{3em}
\newcommand{\N}{\mathbb{N}}
\newcommand{\Z}{\mathbb{Z}}
\newcommand{\Q}{\mathbb{Q}}
\newcommand{\R}{\mathbb{R}}
\newcommand{\C}{\mathbb{C}}
\newcommand{\F}{\mathbb{F}}
\newcommand{\A}{\mathbb{A}}
\newcommand{\PP}{\mathbb{P}}
\newcommand{\E}{\mathbb{E}}
\newcommand{\eps}{\varepsilon}
\newcommand{\dd}{\,\mathrm{d}}
\begin{document}
\section*{UnsolvedMath ID: 800148}
\textbf{Source problem code:} catalogue entry
\subsection{Definitions and mathematical statement}
For the conventional interpretation as races between primes in arithmetic progressions, let \(q\ge 2\) be an integer and let \(a_1,\ldots,a_r\) be pairwise distinct reduced residue classes modulo \(q\), where \(3\le r\le\varphi(q)\) and \(\varphi(q)=\#(\Z/q\Z)^\times\). Define
\[
\pi(x;q,a)=\#\{p\le x:p\text{ is prime and }p\equiv a\pmod q\}
\]
and the ordering event
\[
R(q;a_1,\ldots,a_r)=\{x\ge2:\pi(x;q,a_1)>\cdots>\pi(x;q,a_r)\}.
\]
When the limit exists, define its logarithmic density by
\[
\delta(q;a_1,\ldots,a_r)=\lim_{X\to\infty}\frac{1}{\log X}\int_2^X\mathbf 1_{R(q;a_1,\ldots,a_r)}(x)\,\frac{\mathrm dx}{x}.
\]
Here \(\mathbf 1_R\) denotes the indicator of \(R\). This is an explicit conventional framework; the supplied excerpt does not itself define the races, the density convention, or hypotheses ensuring existence of these limits.

\textbf{Requested task.} Numerically compute further three-way races, or races involving more competitors. In particular, determine, for example, the \(100\) most biased three-way races. The reference value for a three-way ordering density is \(1/6\). The excerpt does not specify whether ``most biased'' means the largest ordering densities, the largest absolute deviations from \(1/6\), or another statistic. It also does not specify the admissible moduli, restrictions on the competitors, or equivalences under which two races count as the same. Consequently, a unique ranked list cannot be specified from this excerpt alone; these choices must be supplied before a precise exhaustive-ranking task is fixed.

The source additionally highlights the potentially harder task of bounding three-way densities sufficiently sharply to certify that such a list is exhaustive. It reports that Fiorilli and Martin [FiM] determined the \(117\) most biased two-way races. It attributes to Lamzouri [Lam] a convergence rate \(O(1/\log q)\) toward \(1/6\) for three-way race densities, compared with a rate \(q^{-1/2+o(1)}\) toward \(1/2\) for two-way densities. These are supplied background claims: their hypotheses, ranges of uniformity, and precise bibliographic references are absent from the record. They are not asserted here as unconditional statements for every tuple defined above.
\subsection{Short English statement}
Compute more races among primes in three or more residue classes, and identify, for example, the \(100\) most biased three-way races. The source also raises the challenge of proving that the resulting list is exhaustive. The meaning of bias, the search domain, the treatment of equivalent races, and the assumptions needed for the densities remain unspecified.
\subsection{Sources}
\begin{itemize}
\item Local UnsolvedMath registry, permanent ID \(800148\), exact source code \texttt{catalogue entry}; title: \emph{Comparative Prime Number Theory --- Problem 2}. Supplied upstream revision: \texttt{d371ee917ac9647fbb06cd6fc529deec6f3dda82}. The registry labels this record \emph{open}; this is supplied curation, not an independently verified status assessment.
\item Hamieh, Kadiri, Martin, and Ng, \emph{Comparative Prime Number Theory Problem List} (2024), \url{https://arxiv.org/abs/2407.03530}. The supplied provenance identifies the source item as ``Problem 2 (\#17)'', records access on 2026-10-16, and identifies the extraction method as \texttt{arxiv\_latex}. These provenance details are supplied, not independently verified here. The extracted statement instead begins ``Problem 1''.
\item The excerpt attributes the proposed task to Greg Martin, University of British Columbia. The registry's \texttt{proposed\_by} field lists Hamieh, Kadiri, Martin, and Ng, with proposed year 2024.
\item [FiM]: Fiorilli and Martin, cited in the supplied excerpt for determining the \(117\) most biased two-way races. No title, publication details, or separate link for [FiM] is supplied.
\item [Lam]: Lamzouri, cited in the supplied excerpt for the comparison of three-way and two-way density convergence rates. No title, publication details, or separate link for [Lam] is supplied.
\end{itemize}
\end{document}
